% Простіший чорний текст розділів 1--3; основа — узгоджена parallel-версія.
\begingroup\color{black}
\hypersetup{citecolor=black,linkcolor=black}

\section{Вступ}
\label{sec:introduction}

Епідемія поширюється на кількох рівнях. Усередині спільноти зараження може залежати від конкретних контактів між людьми \citep{keeling2005,pastorsatorras2015,kiss2017}. Для зв'язків між спільнотами дані часто менш докладні: наприклад, відома лише кількість пасажирів, обсяг поїздок на роботу або загальна вага зв'язку між двома місцями \citep{balcan2009,barbosa2018human}. Тому потрібна модель, яка поєднує докладну контактну мережу всередині кожної спільноти із загальними даними про зв'язки між спільнотами.

Така ситуація виникає на практиці. Локальну контактну мережу можна спостерігати, відновлювати або генерувати за відомими характеристиками, зокрема кластеризацією та середньою довжиною найкоротших шляхів \citep{eubank2004,danon2011}. Натомість зв'язки між спільнотами можуть бути відомі лише з транспортних, міграційних даних, даних про поїздки на роботу чи даних мобільних операторів \citep{barbosa2018human,chang2021}. Вони дають загальну силу зв'язку, але не визначають людей, через яких інфекція може перейти між спільнотами. Отже, ми можемо знати вагу зв'язку, не знаючи конкретних вузлів-мостів та їхнього стану.

Якщо відомі всі контакти, епідемію можна описати неперервним марковським ланцюгом (CTMC) і точно моделювати методом Gillespie \citep{gillespie1977exact,gibson2000}. Є також ефективні методи для часових, великих і розріджених мереж \citep{vestergaard2015,kuryliak2025simulating}. Проте їм усе одно потрібні конкретні контакти, здатні передавати інфекцію. Метапопуляційні моделі використовують загальні потоки між спільнотами, але зазвичай замінюють і докладну локальну мережу кількостями людей у різних станах \citep{colizza2007reaction,colizza2008,balcan2010modeling,ajelli2010}. Наша задача лежить між цими двома випадками: зберегти випадкове поширення між окремими вузлами всередині спільнот, використовуючи лише загальні дані для зв'язків між ними.

Базовий двомасштабний симулятор представлено в попередній конференційній роботі \citep{kuryliak2025towards}. Для кожної спільноти він запускає окремий локальний симулятор. Макрорівень використовує поточні стани спільнот і ваги зв'язків, щоб генерувати зараження, які надходять з іншої спільноти. У цій статті ми докладно описуємо метод і досліджуємо його похибку. Зокрема, з'ясовуємо, що втрачається, коли справжні міжспільнотні ребра та їхні кінцеві вузли замінено загальною вагою зв'язку, а імпортоване зараження потрапляє на вибраний сприйнятливий вузол.

Головне питання: коли такого агрегованого зв'язку достатньо для відтворення потрібної нам характеристики повномережевого CTMC? Ми досліджуємо часи прибуття інфекції, їхні розподіли та можливість передавання через справжні контакти. Близькі часи прибуття ще не гарантують фізично можливого передавання: агрегована модель може заразити іншу спільноту до того, як заразиться будь-який справжній вузол-мост у спільноті-джерелі. У контрольованих експериментах ми знаємо повну контактну мережу й використовуємо її як еталон. Двомасштабний симулятор отримує лише локальні мережі та загальні ваги зв'язків. Це перевірка методу, а не калібрування моделі конкретної епідемії.

Стаття має чотири основні результати.
\begin{enumerate}
\item Ми повністю описуємо двомасштабний симулятор, зокрема узгодження локальних симуляторів за часовою сіткою, використане в експериментах (Розділ~\ref{sec:method}, Алгоритм~\ref{alg:scheduler}).
\item Ми показуємо, що саме моделює ідеальний метод із неперервним часом. Він точний для \emph{сурогатної мережі}: справжні міжспільнотні ребра замінено слабкими зв'язками між усіма парами вузлів зі збереженням загальної ваги (Твердження~\ref{prop:surrogate}). Загалом це не є точним моделюванням початкової фізичної мережі. Далі ми визначаємо джерела похибки та аналізуємо один експорт інфекції (Твердження~\ref{prop:ordering}).
\item Ми порівнюємо метод із повномережевим моделюванням Gillespie для одиничних експортів і ланцюгів спільнот. Перестановочні перевірки показують, які відмінності можуть виникнути лише через скінченну кількість запусків (Розділи~\ref{sec:design} і~\ref{sec:results}).
\item Ми пропонуємо емпіричний показник точності для наших тестів. Він порівнює розкид часів зараження всередині спільноти з типовим часом передавання через одне міжспільнотне ребро (Розділ~\ref{sec:res_descriptor}).
\end{enumerate}

Ми розглядаємо динаміку «сприйнятливий--інфікований» (SI) на статичних синтетичних мережах і використовуємо один фіксований крок синхронізації. Розділ~\ref{sec:related} пояснює зв'язок із попередніми методами. Розділ~\ref{sec:method} описує симулятор. У Розділі~\ref{sec:analysis} пояснено, яку інформацію агрегований зв'язок зберігає та втрачає. Розділи~\ref{sec:design} і~\ref{sec:results} містять експерименти та результати. У Розділі~\ref{sec:discussion} обговорено застосування, обчислювальні витрати й обмеження.

\clearpage
\section{Пов'язані праці}
\label{sec:related}

\paragraph{Моделювання випадкових подій.}
Прямий метод Gillespie вибирає наступну подію та її час за поточними інтенсивностями подій \citep{gillespie1977exact}. Метод наступної реакції підтримує окремий годинник для кожної можливої події \citep{gibson2000}. Обидва методи точні для CTMC, який вони моделюють. Якщо інтенсивність змінюється з часом, її можна накопичувати до досягнення випадкового порога \citep{ethier1986}. \citet{anderson2007} використав цю ідею для систем зі змінними інтенсивностями та затримками. Наш макрогодинник працює за тією самою ідеєю, але його інтенсивність змінюється разом зі станами локальних епідемій. Мережеві симулятори можуть зменшувати обсяг роботи, оновлюючи лише інтенсивності, на які вплинула остання подія \citep{vestergaard2015,kuryliak2025simulating}. Наші локальні симулятори використовують реалізацію з \citet{kuryliak2025simulating}.

Деякі багатомасштабні методи прискорюють моделювання хімічних реакцій, наближено описуючи реакції, що відбуваються значно частіше за інші \citep{haseltine2002,cao2005slow}. Наш метод прибирає іншу інформацію. Він зберігає локальні зараження, але не знає, які саме міжспільнотні ребра зараз здатні передавати інфекцію. Попри цю відмінність, співвідношення швидкостей локального поширення та міжспільнотного передавання допомагає пояснити похибку.

\paragraph{Моделі з кількома спільнотами.}
Метапопуляційні моделі описують кожну підпопуляцію кількостями людей у різних станах і зв'язують підпопуляції потоками мобільності. Їх використовували для дослідження глобального поширення та умов проникнення інфекції \citep{colizza2007reaction,colizza2008,balcan2009,balcan2010modeling}. Докладні дані мобільності також можна поєднувати з локальними моделями передавання \citep{chang2021}. Агентні моделі, своєю чергою, представляють окремих людей у всій системі \citep{eubank2004,kerr2021}. \citet{ajelli2010} порівняли ці два підходи. Досліджуваний тут метод зберігає окремі вузли та контакти всередині спільноти, але замінює невідомі індивідуальні контакти між спільнотами одним агрегованим стохастичним зв'язком.

\paragraph{Два рівні перемішування та спільне моделювання.}
Деякі епідемічні моделі поєднують контакти в малих локальних групах із рівномірним перемішуванням на ширшому рівні \citep{ball1997}. Наша сурогатна мережа має подібну ідею, але рівномірне перемішування діє лише між зв'язаними парами спільнот, а його силу визначає загальна вага зв'язку.

Реалізація також має риси спільного моделювання (co-simulation). У ньому окремі симулятори підсистем рухаються вперед за спільною часовою сіткою та обмінюються обмеженим набором даних \citep{gomes2018cosim}. У нас цими підсистемами є локальні симулятори спільнот. Вони передають макрорівню кількості інфікованих і сприйнятливих вузлів. Макрорівень використовує ці дані, щоб згенерувати час міжспільнотного зараження та вибрати спільноту-одержувача й цільовий вузол. Отже, локальні симулятори можуть залишатися «чорними скриньками» для планувальника. У цій реалізації дані, якими вони обмінюються, використовуються для створення випадкових подій зараження.

\clearpage
\section{Метод двомасштабного моделювання}
\label{sec:method}

В експериментах використано одну реалізацію двомасштабного симулятора, яка узгоджує спільноти за часовою сіткою. Щоб зрозуміти її похибку, ми також розглядаємо процес із тими самими інтенсивностями, але з точним визначенням часу кожної події. Цей теоретичний опис називаємо \emph{ідеальним процесом}. Він дає змогу відокремити похибку заміни мережі від додаткової похибки часової сітки. Ідеальний процес точний для сурогатної мережі, визначеної далі; це не означає точності для початкової фізичної мережі. Нижче пояснимо еталонну модель, доступні дані та кроки реалізованого симулятора.

\subsection{Повномережева еталонна модель}
\label{sec:reference}

Нехай $G=(V,E)$ --- контактна мережа. Її вузли поділено на $M$ спільнот, які не перетинаються:
\begin{equation}
\label{eq:partition}
V=V_1\cup V_2\cup\cdots\cup V_M,
\qquad N_i=|V_i|.
\end{equation}
Тут $V_i$ --- множина вузлів спільноти $i$, а $N_i$ --- кількість її вузлів. Позначаємо через $E_{ii}$ контакти всередині спільноти $i$, а через $E_{ij}$ --- контакти від вузлів спільноти $i$ до вузлів спільноти $j$. Контакт представляємо як напрямлений канал передавання. Тому ненапрямлений контакт дає два канали --- по одному в кожному напрямку.

Кожен вузол є сприйнятливим ($S$), інфікованим ($I$) або видаленим ($R$). Повний стан мережі в час $t$ позначаємо $X(t)$, а стан спільноти $i$ --- $X_i(t)$. Множини сприйнятливих та інфікованих вузлів цієї спільноти позначаємо $S_i(t)$ та $I_i(t)$. Їхні кількості --- $s_i(t)$ та $y_i(t)$, а частки --- $\sbar_i(t)=s_i(t)/N_i$ та $\ibar_i(t)=y_i(t)/N_i$. Далі $x$ означає конкретний стан мережі; $x^{v:I}$ --- той самий стан після зараження вузла $v$, а $x^{u:R}$ --- після видалення вузла $u$.

Повна модель має два прості правила. Сприйнятливий вузол $v$ отримує інфекцію від кожного інфікованого сусіда $u$ з інтенсивністю $\beta_{uv}$. Інфікований вузол видаляється з інтенсивністю $\gamma$. Для вузлів у відповідних станах:
\begin{equation}
\label{eq:full_generator}
q_v(x)=\sum_{\substack{u:(u,v)\in E\\x_u=I}}\beta_{uv}
\quad (x_v=S),
\qquad r_u(x)=\gamma\quad (x_u=I).
\end{equation}
Тут $q_v$ --- сумарна інтенсивність зараження вузла $v$, а $r_u$ --- інтенсивність видалення вузла $u$. Для несприйнятливого вузла $q_v=0$; для неінфікованого вузла $r_u=0$. Інтенсивність визначає частоту події: наприклад, за сталої інтенсивності $q$ середній час очікування дорівнює $1/q$. Ці правила також визначають генератор CTMC $\mathcal{L}_G$, використаний далі; ця назва не додає нових правил епідемії. Для міжспільнотного каналу записуємо $\beta_{uv}=\beta_{\mathrm{out}}w_{uv}$, де $w_{uv}$ --- вага контакту. Точне моделювання цих повномережевих інтенсивностей позначаємо \Full{}. Усі експерименти в статті мають динаміку SI, тобто $\gamma=0$.

\subsection{Агрегований інтерфейс}
\label{sec:interface}

Двомасштабний симулятор \TS{} знає кожен локальний граф $G_i=(V_i,E_{ii})$ та його інтенсивності. Він не знає конкретних міжспільнотних каналів у $E_{ij}$. Для кожної напрямленої пари спільнот йому відома лише загальна вага:
\begin{equation}
\label{eq:weight}
W_{ij}=\sum_{(u,v)\in E_{ij}}w_{uv},\qquad i\neq j.
\end{equation}
Якщо всі ваги дорівнюють одиниці, $W_{ij}$ --- кількість каналів від $V_i$ до $V_j$. У прикладній задачі цю величину можна оцінити за транспортними, міграційними даними, даними мобільності або контактів. Макрорівень також отримує поточні кількості сприйнятливих та інфікованих вузлів від кожного локального симулятора. Він не знає, які вузли є кінцями справжніх міжспільнотних ребер. Основні відмінності наведено в Таблиці~\ref{tab:information}; Рисунок~\ref{fig:schematic} показує їх на схемі.

\begin{table}[tbp]
\color{black}\centering
\caption{\textcolor{black}{Дані та правила двох моделей. В усіх експериментах \TS{} вибирає цільовий вузол рівномірно серед сприйнятливих вузлів.}}
\label{tab:information}
\small
\begin{tabular}{>{\raggedright\arraybackslash}p{0.42\linewidth}>{\raggedright\arraybackslash}p{0.23\linewidth}>{\raggedright\arraybackslash}p{0.25\linewidth}}
\toprule
Дані або дія & \Full{} & \TS{}\\
\midrule
Локальні контакти, інтенсивності та стани вузлів & відомі & відомі локальним симуляторам\\
Конкретні міжспільнотні ребра та вузли-мости & відомі & невідомі\\
Загальна міжспільнотна вага $W_{ij}$ & отримана з ребер & задана безпосередньо\\
Інтенсивність міжспільнотного зараження & зі станів кінців справжніх ребер & з $\ibar_i$, $\sbar_j$ та $W_{ij}$\\
Цільовий вузол зараження & кінець ребра, яке передало інфекцію & сприйнятливий вузол, вибраний за правилом\\
Кількість можливих SI-заражень через набір каналів & не більше кількості каналів & не більше $s_j$ цільових заражень\\
\bottomrule
\end{tabular}
\end{table}

\begin{figure}[tbp]
\color{black}\centering
\begin{tikzpicture}[
  x=1cm,y=1cm,
  font=\scriptsize,
  >=Latex,
  person/.style={circle,draw=black!60,fill=blue!18,minimum size=4.2mm,inner sep=0pt},
  infected/.style={person,fill=orange!70},
  bridge/.style={person,very thick,draw=black},
  bridgeinf/.style={infected,very thick,draw=black},
  commbox/.style={draw=black!45,dashed,rounded corners=2pt,inner sep=3pt},
  macrobox/.style={draw=black!70,rounded corners=2pt,fill=black!5,align=center,inner sep=3pt},
  link/.style={draw=black!55},
  weak/.style={draw=black!24,line width=.25pt},
  flow/.style={-{Latex[length=1.7mm]},draw=black!65}
]
% (a) Full microscopic interface
\node[font=\bfseries] at (0,2.25) {(a) Full network};
\node[infected] (a1) at (-1.65,1.20) {};
\node[person]   (a2) at (-1.55,0.35) {};
\node[bridgeinf](a3) at (-0.75,0.82) {};
\node[person]   (a4) at (-0.75,1.55) {};
\node[bridge]   (b1) at (0.65,1.15) {};
\node[person]   (b2) at (1.45,1.55) {};
\node[person]   (b3) at (1.55,0.55) {};
\node[bridge]   (b4) at (0.70,0.30) {};
\draw[link] (a1)--(a2)--(a3)--(a1) (a1)--(a4)--(a3);
\draw[link] (b1)--(b2)--(b3)--(b4)--(b1) (b1)--(b3);
\draw[very thick] (a3)--(b1);
\draw[very thick] (a3)--(b4);
\node[commbox,fit=(a1)(a2)(a3)(a4)] {};
\node[commbox,fit=(b1)(b2)(b3)(b4)] {};
\node at (-1.2,-0.10) {$V_i$};
\node at (1.15,-0.10) {$V_j$};

% (b) Two-scale representation
\node[font=\bfseries] at (5.0,2.25) {(b) Two-scale simulator};
\node[macrobox] (macro) at (5.0,1.55) {macro interface\\$W_{ij}$, $(\ibar_i,\sbar_j)$\\$\lambda_{ij}$ + placement};
\node[macrobox,minimum width=2.15cm,minimum height=.70cm] (li) at (3.65,0.35) {local CTMC\\$G_i$};
\node[macrobox,minimum width=2.15cm,minimum height=.70cm] (lj) at (6.35,0.35) {local CTMC\\$G_j$};
\draw[flow] (li.north) -- node[left] {summaries} (macro.south west);
\draw[flow] (macro.south) -- node[right] {import} (lj.north);
\draw[flow] (lj.east) -- ++(0,1.0) -| node[pos=1.08,right] {summaries} (macro.east);
\node at (5.0,-0.35) {node-level dynamics retained locally};

% (c) Surrogate interface
\node[font=\bfseries] at (10.25,2.25) {(c) Surrogate network};
\node[infected] (c1) at (8.85,1.45) {};
\node[person]   (c2) at (8.75,0.55) {};
\node[person]   (c3) at (9.55,1.05) {};
\node[person]   (c4) at (9.45,0.20) {};
\node[person]   (d1) at (11.00,1.45) {};
\node[person]   (d2) at (10.90,0.55) {};
\node[person]   (d3) at (11.70,1.05) {};
\node[person]   (d4) at (11.60,0.20) {};
\foreach \u in {c1,c2,c3,c4}{\foreach \v in {d1,d2,d3,d4}{\draw[weak] (\u)--(\v);}}
\draw[link] (c1)--(c2)--(c4)--(c3)--(c1);
\draw[link] (d1)--(d2)--(d4)--(d3)--(d1);
\node[commbox,fit=(c1)(c2)(c3)(c4)] {};
\node[commbox,fit=(d1)(d2)(d3)(d4)] {};
\node[align=center] at (10.25,-0.35) {$W_{ij}/(N_iN_j)$ per\\cross-community channel};
\end{tikzpicture}
\caption{\textcolor{black}{Три подання зв'язку між спільнотами. (a) У повній мережі інфекція переходить лише через показані фізичні ребра та їхні вузли-мости. (b) Двомасштабний метод зберігає окремий локальний симулятор для кожної спільноти. Макрорівень отримує кількості вузлів у різних станах, використовує $W_{ij}$ для генерування міжспільнотного зараження та вибирає його ціль рівномірно серед сприйнятливих вузлів. (c) Ідеальний двомасштабний процес еквівалентний процесу на сурогатній мережі, де кожна напрямлена пара вузлів між двома спільнотами має слабкий зв'язок. Загальна вага всіх цих зв'язків дорівнює $W_{ij}$. Помаранчеві вузли --- інфіковані, сині --- сприйнятливі.}}
\label{fig:schematic}
\end{figure}

\subsection{Локальні симулятори}
\label{sec:local}

Кожна спільнота має окремий симулятор Gillespie для свого графа $G_i$. Він застосовує інтенсивності зараження та видалення з~\eqref{eq:full_generator}, використовуючи лише локальні ребра $E_{ii}$. Реалізація оновлює лише ті інтенсивності, на які вплинула остання подія, як у \citet{kuryliak2025simulating}. Кожен локальний симулятор має власний потік випадкових чисел.

Макрорівень може просунути локальний симулятор до заданого часу. Симулятор виконує локальні події до цього моменту. Потім макрорівень може імпортувати інфекцію, змінивши стан вибраного сприйнятливого вузла на $I$. Це змінює локальні інтенсивності. Реалізація зберігає невикористану частину випадкового порога наступної локальної події та перераховує час очікування за новою сумарною інтенсивністю. Тип події та її вузол вибираються за оновленими інтенсивностями. Новий локальний поріг генерується лише тоді, коли збереженого порога немає. Кожен локальний годинник має власний поріг, окремий від макропорога, описаного в Розділі~\ref{sec:timing}. За заданих часів імпорту та цільових вузлів локальний симулятор виконує правила CTMC своєї мережі. Наближення виникає через спосіб генерування імпортів і вибору їхніх цільових вузлів.

\subsection{Макрорівень}
\label{sec:macro}

Макрорівень представляє спільноту її поточними частками вузлів, а напрямлений зв'язок --- вагою $W_{ij}$. Інтенсивність зараження зі спільноти $i$ до спільноти $j$:
\begin{equation}
\label{eq:macro_rate}
\lambda_{ij}(t)=\beta_{\mathrm{out}}W_{ij}\ibar_i(t)\sbar_j(t),
\qquad \Theta(t)=\sum_{i\neq j}\lambda_{ij}(t).
\end{equation}
Кожен множник має простий зміст. $\beta_{\mathrm{out}}W_{ij}$ задає загальну силу зв'язку; $\ibar_i$ --- частку інфікованих у спільноті-джерелі; $\sbar_j$ --- частку сприйнятливих у спільноті-одержувачі. Сума $\Theta$ --- інтенсивність будь-якого міжспільнотного зараження.

Коли така подія відбувається, спочатку вибирають спільноту-джерело та спільноту-одержувача $(i,j)$ з імовірністю $\lambda_{ij}/\Theta$. Потім у спільноті $j$ рівномірно вибирають один сприйнятливий вузол і заражають його. Якщо таких вузлів 20, кожен має імовірність $1/20$. Вибір не враховує справжні вузли-мости, оскільки агрегований інтерфейс їх не знає. У подальших формулах цю імовірність позначаємо $\Pi_{ij}(v\mid x)=1/s_j$ для $v\in S_j$, де $x$ --- поточний стан мережі. Отже, $\Pi$ просто позначає імовірність вибору конкретного сприйнятливого вузла. Інші правила можливі, якщо є інформація про вузли-мости. Додатковий калібрувальний множник у реалізації дорівнює одиниці в усіх експериментах.

Ця агрегована інтенсивність загалом наближено описує передавання через справжні фізичні ребра. Для заданих кількостей вузлів у станах вона дорівнює очікуваній інтенсивності за рівномірного випадкового вибору кінців ребер, перш ніж враховано їхні конкретні ідентичності, стани та історію епідемії. Це не гарантує правильної інтенсивності для вже реалізованої епідемії на конкретних ребрах. Для рівномірного зв'язку «усі з усіма» вона є фактичною інтенсивністю, як пояснено в Розділі~\ref{sec:analysis}. Значення $\Theta$ змінюється після локальної події або імпорту, які змінюють кількості вузлів у станах.

\subsection{Час подій та синхронізація}
\label{sec:timing}

Після макроподії в час $t_0$ генеруємо випадковий поріг $E\sim\Exp(1)$: експоненційне випадкове число із середнім один. Починаючи з нуля, накопичуємо загальну інтенсивність $\Theta$. Наступна макроподія відбувається в перший момент $t^\ast$, коли накопичена величина досягає $E$:
\begin{equation}
\label{eq:hazard_integral}
\int_{t_0}^{t^\ast}\Theta(s)\,ds=E.
\end{equation}
Інакше кажучи, площа під кривою інтенсивності зростає до випадкового порога. Якщо поріг не досягнуто, наступної макроподії немає. Це стандартний спосіб визначення часу події за накопиченою інтенсивністю \citep{ethier1986,anderson2007}. За сталої $\Theta$ час очікування дорівнює $E/\Theta$. Якщо згенерувати $E=-\ln U$, де $U$ рівномірно розподілене на $(0,1)$, отримуємо звичний час очікування Gillespie. Коли локальна подія змінює $\Theta$, зберігаємо той самий поріг і продовжуємо накопичення з новою інтенсивністю. Новий макропоріг генеруємо лише після макроподії. Це той самий годинник за накопиченою інтенсивністю, який описано в конференційній роботі. Точне визначення першого досягнення порога задає ідеальний процес; описаний нижче розрахунок на часовій сітці наближено визначає цей момент. У Розділі~\ref{sec:surrogate} показано, що ідеальний процес точний для сурогатної мережі.

Реалізований планувальник узгоджує локальні симулятори на інтервалах довжини $\tau$. У конференційному описі цього годинника також використано оцінку інтенсивності в середині кроку. Тут уточнюємо, як отримати її від локальних симуляторів. Нехай $A$ --- уже накопичена величина від попередньої макроподії. Для поточного інтервалу $[t,t_g)$ копіюємо кожен локальний симулятор разом зі станом генератора випадкових чисел і просуваємо копію до середини інтервалу. За копіями обчислюємо оцінку загальної макроінтенсивності $\Theta_{\mathrm{mid}}$. Оцінений приріст за залишок інтервалу дорівнює $\Theta_{\mathrm{mid}}(t_g-t)$. Якщо $A$ плюс цей приріст менше за $E$ (випадковий макропоріг), справжні симулятори просуваються до $t_g$. Інакше час події оцінюємо як
\begin{equation}
\label{eq:interpolation}
t^\ast=t+\frac{E-A}{\Theta_{\mathrm{mid}}}.
\end{equation}
Тут $E$ --- випадковий макропоріг, а $A$ --- уже накопичена до нього величина. Чисельник $E-A$ --- залишок до порога; ділення на оцінену інтенсивність дає залишковий час очікування. Наприклад, за $E=1$, $A=0.4$ та $\Theta_{\mathrm{mid}}=0.2$ оцінений час очікування дорівнює трьом одиницям часу. Планувальник просуває справжні симулятори до $t^\ast$, за їхніми станами в цей момент вибирає спільноту-джерело та спільноту-одержувача, вибирає сприйнятливий цільовий вузол і застосовує зараження. Повну послідовність наведено в Алгоритмі~\ref{alg:scheduler}.

\begin{algorithm}[H]
\color{black}
\caption{\textcolor{black}{Один запуск двомасштабного симулятора з узгодженням за часовою сіткою.}}
\label{alg:scheduler}
\small
\fbox{\begin{minipage}{0.94\linewidth}
\textbf{Вхід:} локальні графи $G_i$, ваги $W_{ij}$, інтенсивності, крок $\tau$, горизонт $T_{\mathrm{end}}$ та початковий стан.
\begin{enumerate}
\item Установити $t=0$ і $A=0$; згенерувати $E\sim\Exp(1)$.
\item Поки $t<T_{\mathrm{end}}$ й інфекція може поширюватися, встановити $t_g=\min(t+\tau,T_{\mathrm{end}})$ та обробити інтервал до $t_g$.
\item Скопіювати кожен локальний симулятор разом зі станом генератора випадкових чисел. Просунути копії до середини залишку інтервалу та обчислити $\Theta_{\mathrm{mid}}$ за~\eqref{eq:macro_rate}.
\item Якщо $A+\Theta_{\mathrm{mid}}(t_g-t)<E$, просунути справжні симулятори до $t_g$, додати цю величину до $A$, встановити $t=t_g$ і повернутися до кроку 2.
\item Інакше обчислити $t^\ast$ за~\eqref{eq:interpolation} та просунути справжні симулятори до $t^\ast$.
\item Вибрати $(i,j)$ з імовірністю $\lambda_{ij}(t^\ast)/\Theta(t^\ast)$. Рівномірно вибрати $v$ з $S_j(t^\ast)$ та імпортувати зараження у $v$. Оновити локальний симулятор одержувача, зберігаючи невикористаний експоненційний поріг і вибираючи локальну подію за новими інтенсивностями.
\item Установити $t=t^\ast$ і $A=0$, згенерувати нове $E\sim\Exp(1)$ та так само обробити залишок поточного інтервалу.
\end{enumerate}
\end{minipage}}
\end{algorithm}

Копіювання станів генераторів випадкових чисел важливе: якщо до середини інтервалу не відбудеться імпорт, копія відтворить ту саму траєкторію, що й справжній симулятор. Оцінка середньої точки точна на інтервалі без локальних подій. Її похибка визначення часу та використання майбутнього стану не більш як на $\tau/2$ зникають зі зменшенням $\tau$ до нуля.

Часова сітка потрібна, оскільки планувальник працює з локальними симуляторами як із «чорними скриньками». Якщо вони також повідомляють часи наступних локальних подій, планувальник може об'єднати ці часи та точно накопичувати $\Theta$, яка стала між подіями. Такий подієвий варіант усуває чисельну похибку синхронізації. В експериментах цієї статті використано реалізований сітковий варіант із $\tau=1$.

\subsection{Обчислювальна структура та застосування}
\label{sec:cost}

\Full{} зберігає всі локальні та міжспільнотні канали, а обсяг роботи на одну подію зростає зі степенем вузла \citep{kuryliak2025simulating}. \TS{} зберігає локальні графи та матрицю загальних ваг розміру $M\times M$. У конференційній роботі обговорено таке скорочення даних про міжспільнотні контакти \citep{kuryliak2025towards}. Його значення залежить від задачі. Заміна багатьох індивідуальних міжспільнотних ребер однією вагою може зменшити обсяг збережених даних. Якщо таких ребер лише кілька, економія може бути малою; локальні графи все одно потрібно зберігати.

Окремі локальні симулятори також створюють можливість паралельного виконання. Коли планувальник визначив наступний цільовий час ($t_g$ або $t^\ast$), локальні симулятори можна незалежно просунути до цього самого часу. Копії для оцінки середньої точки також можна обробляти паралельно. Потім макрорівень має зібрати кількості вузлів і застосувати імпорт, якщо він запланований, перш ніж продовжувати. У досліджуваній реалізації локальні симулятори просуваються послідовно; паралельне виконання в межах одного запуску є можливим розширенням, а не виміряним прискоренням. Локальні симулятори також можуть мати різні налаштування для різних спільнот.

Сітковий планувальник додає роботу через копіювання локальних станів та узгодження спільнот. Тому ні менше використання пам'яті, ні швидше виконання не гарантовані. Метод корисний, коли конкретні міжспільнотні контакти недоступні або надто численні для явного подання. Виміряний час виконання наведено в Розділі~\ref{sec:disc_cost}; головний предмет цієї статті --- похибка наближення.

\clearpage
\endgroup
